% Libro 2 · anexos: formulario, plan de estudio y simulacro final
\makeatletter
\cleardoublepage\renewcommand{\zuka@mundoactual}{Formulario}
\phantomsection\addcontentsline{toc}{chapter}{Formulario}
{\sffamily\fontsize{26}{30}\selectfont\firalight Formulario\par}\vspace{2pt}
{\color{zlime}\rule{18mm}{1.6pt}}\par\vspace{6pt}

\begin{multicols}{2}\small
\kicker{Orden (1)}
\begin{itemize}
  \item Fila de casillas; fijas y bloques primero.
  \item Mesa redonda: fija a uno; frente a $k$ está $k \pm \tfrac n2$.
  \item Tabla: una ✓ por fila y por columna.
  \item Tiene que: en todas. Podría: en una. No puede: en ninguna.
\end{itemize}
\kicker{Lógica (2)}
\begin{itemize}
  \item $p \to q$ es F solo si $p$ V y $q$ F.
  \item $p \to q \equiv \neg q \to \neg p \equiv \neg p \vee q$.
  \item $\neg(p \wedge q) \equiv \neg p \vee \neg q$; $\neg(p \vee q) \equiv \neg p \wedge \neg q$.
  \item $\neg(\text{todos}) =$ algunos no; $\neg(\text{algunos}) =$ ninguno.
  \item Válidos: ponens, tollens.
\end{itemize}
\kicker{Conjuntos (3)}
\begin{itemize}
  \item $n(A \cup B) = n(A) + n(B) - n(A \cap B)$.
  \item Encuestas: del centro hacia afuera. Subconjuntos: $2^n$.
\end{itemize}
\columnbreak
\kicker{Conteo y probabilidad (4, 5)}
\begin{itemize}
  \item «Y luego»: $\times$. «O»: $+$.
  \item $n!$; $P(n, k) = \tfrac{n!}{(n-k)!}$; $\binom nk = \tfrac{n!}{k!(n-k)!}$.
  \item Repetidos: $\tfrac{n!}{a!b!}$. Circular: $(n - 1)!$.
  \item Caminos $a \times b$: $\binom{a+b}{a}$.
  \item $P = \tfrac{\text{favorables}}{\text{posibles}}$; al menos uno: $1 - P(\text{ninguno})$.
\end{itemize}
\kicker{Patrones y códigos (6, 7)}
\begin{itemize}
  \item Diferencias, cocientes, intercaladas, conocidas.
  \item $a_n = a_1 + (n - 1)d$; $a_n = a_1 r^{n-1}$.
  \item Ciclo $k$: posición $N \equiv N \bmod k$.
  \item A $= 1$, …, Z $= 26$; invertido: $27 - p$.
\end{itemize}
\kicker{Espacio y gráficas (8, 9)}
\begin{itemize}
  \item Derecha: N, E, S, O. Espejo de reloj: $12{:}00 -$ hora.
  \item Opuestas en el desarrollo: una cara en medio. Dado: opuestas suman 7.
  \item Torta: $1\,\% = 3.6^\circ$. Mediana: ordenar primero.
\end{itemize}
\end{multicols}

\clearpage\renewcommand{\zuka@mundoactual}{Plan de estudio}
\phantomsection\addcontentsline{toc}{chapter}{Plan de estudio con la app}
{\sffamily\fontsize{26}{30}\selectfont\firalight Plan de ocho semanas\par}\vspace{2pt}
{\color{zlime}\rule{18mm}{1.6pt}}\par\vspace{6pt}
{\small Igual que en el Libro 1: cada día una lección del libro, su explicación y sus ejercicios en Zuka, y el taller. Los sábados,
la evaluación del mundo. Si llevas los dos libros de razonamiento a la vez, alterna: un día el Libro 1, otro el Libro 2.\par}
\vspace{6pt}
\begin{center}\small
\begin{tabularx}{\linewidth}{@{}>{\sffamily\bfseries\color{zolive}}lX>{\raggedleft\arraybackslash}p{24mm}@{}}
\toprule
\textbf{Semana} & \textbf{Lecciones} & \textbf{Evaluación}\\\midrule
1 & 1.1 Fila · 1.2 Mesa redonda · 1.3 Tablas · 1.4 Condiciones & Mundo 1\\
2 & 2.1 Conectivos · 2.2 De Morgan · 2.3 Condicional · 2.4 Silogismos · 2.5 Verdad y mentira & Mundo 2\\
3 & 3.1 Venn · 3.2 Encuestas · 4.1 Principio multiplicativo & Mundo 3\\
4 & 4.2 Permutaciones · 4.3 Combinaciones · 4.4 Complemento · 4.5 Caminos & Mundo 4\\
5 & 5.1 Laplace · 5.2 Independencia · 6.1 Sucesiones & Mundo 5\\
6 & 6.2 Letras y figuras · 6.3 Ciclos · 7.1 Sustitución · 7.2 Deducir la regla & Mundos 6 y 7\\
7 & 8.1 Rejillas · 8.2 Cubos · 8.3 Desarrollos · 8.4 Espejos & Mundo 8\\
8 & 9.1 Leer gráficas · 9.2 Medidas · 9.3 Concluir · Simulacro final & Mundo 9 y simulacro\\
\bottomrule
\end{tabularx}
\end{center}

\clearpage\renewcommand{\zuka@mundoactual}{Simulacro final}
\setcounter{ej}{0}
\phantomsection\addcontentsline{toc}{chapter}{Simulacro final}
\immediate\write\zuka@sol{\noexpand\solmun{Simulacro final}}
\cabecera{SIMULACRO FINAL · LIBRO 2}{Orden, lógica y azar}{16 preguntas · 25 minutos\\sin calculadora}\vspace{6pt}
\begin{ejercicios}
\item En una fila de 4, Rosa está detrás de Saúl, Tina es la primera y Uriel está justo delante de Rosa. ¿Quién es el segundo?
\opciones{Saúl}{Uriel}{Rosa}{no se sabe}
\sol{A. Tina 1; Rosa detrás de Saúl y Uriel justo delante de Rosa: T, S, U, R. (1.1)}
\item Si $p$ es V y $q$ es F, ¿cuál es falsa?
\opciones{$p \vee q$}{$q \to p$}{$\neg q$}{$p \to q$}
\sol{D. V $\to$ F es F; las otras tres son verdaderas. (2.1)}
\item La negación de «Ningún alumno llegó tarde» es
\opcionesd{Todos llegaron tarde.}{Algún alumno llegó tarde.}{Algún alumno no llegó tarde.}{Ninguno llegó temprano.}
\sol{B. (2.2)}
\item «Si estudio, apruebo». No aprobé. Entonces
\opcionesd{estudié}{no estudié}{no se sabe}{aprobé}
\sol{B. (2.3)}
\item De 30 estudiantes, 18 juegan fútbol, 12 baloncesto y 5 ambos. ¿Cuántos no juegan ninguno?
\opciones{5}{0}{10}{15}
\sol{A. $30 - 25$. (3.2)}
\item ¿Cuántos números de 3 cifras distintas se forman con 1, 2, 3, 4, 5?
\opciones{60}{125}{10}{120}
\sol{A. (4.1)}
\item ¿Cuántos anagramas tiene la palabra NENE?
\opciones{24}{12}{6}{4}
\sol{C. $\tfrac{4!}{2!2!}$. (4.2)}
\item ¿De cuántas maneras se eligen 3 delegados entre 8 personas?
\opciones{56}{336}{24}{512}
\sol{A. (4.3)}
\item Se lanzan dos dados. ¿Probabilidad de que la suma sea 7 u 11?
\opciones{$\tfrac29$}{$\tfrac16$}{$\tfrac{1}{18}$}{$\tfrac14$}
\sol{A. Suma 7: 6 casos; suma 11: 2 casos. $\tfrac{8}{36} = \tfrac29$. (5.1)}
\item Una moneda se lanza 3 veces. $P(\text{al menos una cara})$ es
\opciones{$\tfrac78$}{$\tfrac38$}{$\tfrac12$}{$\tfrac18$}
\sol{A. (5.2)}
\item ¿Qué número sigue? $\;1,\ 3,\ 7,\ 13,\ 21,\ \ldots$
\opciones{31}{29}{30}{33}
\sol{A. Diferencias 2, 4, 6, 8, 10. (6.1)}
\item ¿Qué letra sigue? $\;$ B, D, G, K, …
\opciones{O}{P}{N}{Q}
\sol{B. Saltos 2, 3, 4, 5: K (11) $+ 5 = 16$, P. (6.2)}
\item Si GATO se escribe HBUP, ¿cómo se escribe RANA?
\opciones{SBOB}{QZMZ}{SBNB}{SAOA}
\sol{A. (7.1)}
\item Mirando al este, alguien gira a la derecha tres veces. ¿Hacia dónde mira?
\opciones{norte}{sur}{oeste}{este}
\sol{A. Tres derechas $=$ una izquierda. (8.1)}
\item Un reloj visto en un espejo marca las 8:15. La hora real es
\opciones{3:45}{4:15}{3:15}{4:45}
\sol{A. (8.4)}
\item La media de 2, 4, 4, 5, 10 es
\opciones{5}{4}{4.5}{25}
\sol{A. (9.2)}
\end{ejercicios}
\makeatother
